% ===== PRÉAMBULE CANONIQUE — à coller VERBATIM en tête de CHAQUE .tex =====
\documentclass[11pt,a4paper]{article}
\usepackage[utf8]{inputenc}\usepackage[T1]{fontenc}\usepackage{lmodern}
\usepackage[french]{babel}
\usepackage{amsmath,amssymb,mathtools}\usepackage{icomma}
\usepackage[version=4]{mhchem}          % \ce{...} : équations & formules
\usepackage{siunitx}\sisetup{locale=FR} % \qty{}{}, \si{}, \ang{}
\usepackage{chemfig}                    % \chemfig{...} : structures développées
\usepackage{xcolor}\usepackage{colortbl}
\usepackage{tikz}\usepackage{pgfplots}\pgfplotsset{compat=1.18}
\usepackage[most]{tcolorbox}\usepackage{enumitem}\usepackage{array,booktabs}
\usepackage[a4paper,margin=2.2cm]{geometry}
\usetikzlibrary{arrows.meta,calc,angles,quotes,positioning,patterns,backgrounds,shapes.geometric}
\definecolor{primary}{RGB}{222,49,99}\definecolor{primarydark}{RGB}{150,22,55}
\definecolor{defcol}{RGB}{0,114,178}\definecolor{methcol}{RGB}{52,92,148}
\definecolor{excol}{RGB}{0,135,90}\definecolor{attncol}{RGB}{200,40,55}
\tcbset{boxbase/.style={enhanced,breakable,boxrule=0.9pt,arc=2.5pt,
  left=3mm,right=3mm,top=2.3mm,bottom=2.3mm,fonttitle=\bfseries,coltitle=white}}
% --- boîtes TUTORIEL (noms EXACTS, ne pas renommer) ---
\newtcolorbox{cle}[1][]{boxbase,colback=defcol!6,colframe=defcol,
  title={\textsc{À retenir}\ifx\relax#1\relax\relax\else\textmd{ : \itshape #1}\fi}}
\newtcolorbox{methode}[1][]{boxbase,colback=methcol!7,colframe=methcol,sharp corners,
  title={\textsc{Méthode}\ifx\relax#1\relax\relax\else\textmd{ --- \itshape #1}\fi}}
\newtcolorbox{exemplebox}[1][]{boxbase,colback=excol!5,colframe=excol,
  title={\textsc{Exemple}\ifx\relax#1\relax\relax\else\textmd{ : \itshape #1}\fi}}
\newtcolorbox{piege}[1][]{boxbase,colback=attncol!6,colframe=attncol,
  title={\textsc{Piège}\ifx\relax#1\relax\relax\else\textmd{ : \itshape #1}\fi}}
\newtcolorbox{remarque}[1][]{boxbase,colback=black!4,colframe=black!45,
  title={\textsc{Remarque}\ifx\relax#1\relax\relax\else\textmd{ : \itshape #1}\fi}}
% --- boîtes EXERCICE / CORRIGÉ (noms EXACTS, auto-comptées) ---
\newtcolorbox[auto counter]{exo}[1][]{boxbase,colback=primary!4,colframe=primary,
  title={\textsc{Exercice}~\thetcbcounter\ifx\relax#1\relax\relax\else\textmd{ --- \itshape #1}\fi}}
\newtcolorbox[auto counter]{corrige}[1][]{boxbase,colback=excol!4,colframe=excol,
  title={\textsc{Corrigé}~\thetcbcounter\ifx\relax#1\relax\relax\else\textmd{ --- \itshape #1}\fi}}
\newcommand{\R}{\mathbb{R}}\newcommand{\N}{\mathbb{N}}\newcommand{\Z}{\mathbb{Z}}\newcommand{\ds}{\displaystyle}
% (un \usepackage{fancyhdr}+en-tête est possible ; il est ignoré côté web)
% =========================================================================
\begin{document}
\begin{exo}[trouver l'autre isotope]
Le gallium ($Z = 31$) possède deux isotopes naturels. L'un est $\ce{^{69}Ga}$, d'abondance $60\,\%$.
La masse atomique relative moyenne du gallium vaut $\overline{A} = 69{,}8$. Détermine le nombre de
masse $A$ de l'autre isotope, puis écris sa notation.
\end{exo}

\begin{exo}[retrouver les deux abondances]
Le brome possède deux isotopes, $\ce{^{79}Br}$ et $\ce{^{81}Br}$. Sa masse atomique relative moyenne
vaut $\overline{A} = 79{,}9$. Détermine l'abondance de chacun des deux isotopes.
\end{exo}

\begin{exo}[mise en situation : la fluorine $\ce{CaF2}$]
La fluorine est un solide ionique de formule $\ce{CaF2}$ (ions $\ce{Ca^{2+}}$ et $\ce{F^{-}}$). On la
suppose formée uniquement des isotopes $\ce{^{40}_{20}Ca}$ et $\ce{^{19}_{9}F}$.
\begin{enumerate}[label=\alph*)]
  \item Combien y a-t-il de protons au total dans un motif $\ce{CaF2}$ ?
  \item Combien de neutrons au total ?
  \item Calcule le rapport $\dfrac{\text{nombre total de protons}}{\text{nombre total de neutrons}}$.
\end{enumerate}
\end{exo}

\begin{exo}[de l'uma à la mole]
On donne $\SI{1}{uma} \approx \SI{1,66e-27}{\kilo\gram}$ et $N_{\!A} \approx \SI{6,02e23}{\per\mole}$.
\begin{enumerate}[label=\alph*)]
  \item Estime la masse, en kilogrammes, d'un atome de $\ce{^{24}Mg}$ ($24$ nucléons $\approx 24$ uma).
  \item Exprime $\SI{1}{uma}$ en grammes.
  \item Déduis-en la masse, en grammes, d'une mole d'atomes de $\ce{^{12}C}$. Que remarques-tu ?
\end{enumerate}
\end{exo}

\begin{exo}[vrai ou faux — synthèse]
Pour chaque proposition, indique si elle est \textbf{vraie} ou \textbf{fausse}, en justifiant.
\begin{enumerate}[label=\alph*)]
  \item Deux isotopes d'un même élément ont, à l'état neutre, le même nombre d'électrons.
  \item La masse atomique relative moyenne d'un élément est nécessairement un nombre entier.
  \item Une abondance de $0{,}33\,\%$ correspond à la fraction $p = 0{,}0033$.
  \item Entre deux isotopes, celui qui a le plus grand nombre de masse possède davantage de protons.
\end{enumerate}
\end{exo}

\end{document}
